\section{L2, L3 y L4: los límites de responsabilidad importan más que el número de nivel}

La sección anterior explicó el entrenamiento de capacidades; esta sección examina los niveles de conducción autónoma. Lang Xianpeng (郎咸朋) enfatiza que «L3 no es una extensión de L2, sino el precursor de L4». Lo decisivo de este juicio no está en el número, sino en la responsabilidad: L2 es asistencia a la conducción, y el conductor humano es responsable; en L3, bajo condiciones específicas, el sistema empieza a asumir la responsabilidad y avisa a la persona con anticipación cuando necesita que tome el control; L4, en cambio, es aquel en que el sistema es plenamente responsable dentro de una zona delimitada o de escenarios delimitados.

\begin{figure}[H]
\centering
\includegraphics[width=0.92\textwidth]{l3-l4-ladder.png}
\caption{L3 no es una extensión de L2: L3 se parece más a un precursor de L4 que a una mejora lineal de la asistencia a la conducción. Diagrama conceptual propio, elaborado a partir de la conversación entre 00:50:50 y 01:15:15.}
\end{figure}

\begin{knowledgebox}{Lectura de la figura: detrás del cambio de nivel hay un cambio de responsabilidad}
A la izquierda, en L2, la persona es responsable y el sistema asiste; a la derecha, L3/L4 empiezan a exigir que el sistema asuma la responsabilidad y que se definan con claridad el ODD y los límites de la toma de control. Al leer los niveles, no basta con fijarse en cuántas funciones hay: hay que considerar la responsabilidad en caso de accidente, el mecanismo de toma de control, la redundancia de seguridad y los límites de operación.
\end{knowledgebox}

\subsection{Por qué L2 no puede convertirse de forma natural en L3}

Esta sección explica un error de concepto. Muchas personas imaginan L2, L3 y L4 como niveles continuos: cada vez más funciones y una experiencia cada vez mejor. Pero, desde el punto de vista de la responsabilidad, el paso de L2 a L3 es un cambio cualitativo. L2 puede exigir que la persona esté atenta en todo momento y tome el control; L3, en cambio, debe garantizar seguridad suficiente durante el periodo en que el sistema es responsable y, al salir del ODD o de los límites de su capacidad, ofrecer una ventana de toma de control que la persona pueda ejecutar. Esto involucra la regulación, la definición del producto, la estrategia de seguridad y la educación de los usuarios.

\begin{warningbox}{Cuanto mejor es L2, más fácil es que el usuario lo use mal}
Si la experiencia de L2 se acerca a la conducción autónoma, pero la responsabilidad legal y de seguridad sigue recayendo en la persona, el usuario puede confiar en exceso en el sistema. Este es el punto más peligroso del límite entre L2 y L3: la experiencia parece de conducción autónoma, pero la responsabilidad sigue siendo de la persona.
\end{warningbox}

\subsection{ODD y toma de control: cómo se traslada el límite de responsabilidad al producto}

Esta sección lleva la cuestión de los niveles al diseño del producto. L3/L4 no consiste en anunciar un nivel más en una presentación, sino en que el sistema sepa dónde puede asumir la responsabilidad, cuándo debe salir, cómo avisar a la persona para que tome el control y cómo pasar a un estado de riesgo mínimo si la toma de control falla. Por ejemplo, las autopistas, las vías rápidas urbanas, el seguimiento de vehículos en congestión, el estacionamiento, ciertas condiciones climáticas y ciertos rangos de velocidad pueden constituir ODD distintos. En cuanto el sistema se compromete a ser responsable dentro de un ODD, necesita el diseño correspondiente de redundancia de percepción, diagnóstico de fallas, monitoreo de estado y responsabilidad regulatoria.

\begin{importantbox}{Tres preguntas sobre el límite de responsabilidad}
Primera: ¿cuándo puede el sistema asumir la responsabilidad de la conducción? Segunda: ¿cuándo debe el sistema devolver la responsabilidad a la persona o pasar a un estado degradado? Tercera: si la persona no toma el control a tiempo, ¿cómo detiene el sistema el vehículo de forma segura o mantiene el riesgo al mínimo? Mientras estas tres preguntas no tengan una respuesta clara, L3 es solo un término de mercadotecnia.
\end{importantbox}

\subsection{Por qué L3 es el precursor de L4}

Esta sección explica el sentido de «precursor». L3 exige que el sistema sea realmente responsable bajo ciertas condiciones, y eso obliga a la empresa a construir las capacidades que necesita L4: definición del ODD, redundancia de seguridad, estrategia de toma de control, minimización del riesgo, monitoreo del sistema, responsabilidad en caso de accidente y un sistema de validación. Por eso no es solo una mejora del paquete de funciones de L2, sino el campo de entrenamiento para entrar en la etapa en que «el sistema es responsable».

\noindent\begin{tabularx}{\textwidth}{p{0.16\textwidth}p{0.36\textwidth}X}
\toprule
Nivel & Responsabilidad central & Implicación de ingeniería \\
\midrule
L2 & La persona es responsable; el sistema asiste & La experiencia funcional puede ser muy buena, pero la persona debe supervisar. \\
L3 & El sistema es responsable bajo condiciones específicas; requiere un mecanismo de toma de control & Se vuelven decisivos los límites de responsabilidad, la estrategia de salida y la certificación regulatoria. \\
L4 & El sistema es plenamente responsable en zonas o escenarios delimitados & Requiere un ciclo cerrado de seguridad completo y la gestión del dominio de operación. \\
L5 & Conducción autónoma sin restricciones & Por ahora se parece más a un concepto de largo plazo. \\
\bottomrule
\end{tabularx}

\subsection{Resumen de la sección}

Los niveles de conducción autónoma no son números de mercadotecnia, sino una estructura de responsabilidad. El valor principal de L3 es que obliga a la empresa a enfrentar las cuestiones de seguridad, regulación, toma de control y responsabilidad del sistema que requiere L4.
